\begin{flushright}
  \fechaclase{07 de septiembre de 2026}
\end{flushright}

\subsection{Principio de invariancia de La Salle}

\paragraph{Teorema}
Sea $\Omega\subset D$ un conjunto invariante y
$V\colon D\to\mathbb{R}$ una función continuamente diferenciable tal que
\[
  \dot{V}\leq 0 \qquad \text{en }\Omega
\]
Sea $E$ el conjunto de todos los puntos $\Omega$ donde $\dot{V}(X)=0$. Sea
$M$ el mayor conjunto invariante contenido en $E$ entonces toda solución que
inicia en $\Omega$ se aproxima a $M$ conforme $t\to\infty$.

\paragraph{Corolario}
Sea $X=0$ un punto de equilibrio. Sea
$V\colon\mathbb{R}^{n}\to\mathbb{R}$ una función continuamente diferenciable,
radialmente no acotada (crece indefinidamente) y definida positiva, tal que
\[
  \dot{V}(X)\leq 0 \qquad \text{en }\mathbb{R}^{n}
\]
\[
  S=\left\{X\in\mathbb{R}^{n}\mid \dot{V}(X)=0\right\}
\]
y supongamos que ninguna solución, excepto la trivial $X(t)=0$ puede
permanecer idéntica en $S$, entonces el sistema es globalmente
asintóticamente estable.

\paragraph{Corolario}
Sea $X=0$ un punto de equilibrio y $V\colon D\to\mathbb{R}$ una función
definida positiva y continuamente diferenciable en $D$ tal que
\[
  \dot{V}(X)\leq 0
\]
y sea
\[
  S=\left\{X\in D\mid \dot{V}(X)=0\right\}
\]
y solo la solución trivial $X=0$ permanece idéntica en $S$, entonces el
origen es asintóticamente estable.

\subsubsection{Análisis de estabilidad para sistemas lineales}

Sea el sistema LTI
\[
  \dot{X}=AX
\]
que es estable si
\[
  \lambda_i(A)<0
\]
alternativamente se propone la función de Lyapunov
\[
  V=X^{T}PX, \qquad P>0
\]
con derivada
\begin{align*}
  \dot{V}
  &=X^{T}P\dot{X}+\dot{X}^{T}PX\\
  &=X^{T}PAX+(AX)^{T}PX\\
  &=X^{T}PAX+X^{T}A^{T}PX\\
  &=X^{T}(PA+A^{T}P)X
\end{align*}
para que el sistema sea estable
\[
  PA+A^{T}P<0
\]
mediante el complemento de Schur se puede expresar la LMI
\[
  \begin{bmatrix}
    PA+A^{T}P & 0\\
    0 & -I
  \end{bmatrix}<0
\]

Si el sistema no es autónomo
\[
  \dot{X}=AX+Bu
\]
y la derivada de la función es
\begin{align*}
  \dot{V}
  &=X^{T}P(AX+Bu)+(AX+Bu)^{T}PX\\
  &=X^{T}PAX+X^{T}PBu+X^{T}A^{T}PX+(Bu)^{T}PX
\end{align*}

Elegimos
\[
  u=-KX
\]
Así,
\[
  \dot{V}
  =X^{T}PAX+X^{T}A^{T}PX
  -X^{T}PBKX-X^{T}K^{T}B^{T}PX
\]

Debido a que $P$ y $K$ son desconocidas, la LMI se convierte en una BMI, por
lo que hay que hacer un cambio de variable
\[
  F=KP^{-1}
\]
\[
  K=FP
\]

Sustituyendo
\[
  X^{T}PBKX=X^{T}P(BF)PX
\]
\[
  X^{T}K^{T}B^{T}PX=X^{T}P(F^{T}B^{T})PX
\]
así
\[
  \dot{V}
  =X^{T}PAX+X^{T}A^{T}PX
  -X^{T}PBFPX-X^{T}PF^{T}B^{T}PX
\]
